\clearpage
\VerdierPhysicalPageAnchor{0268}%
\begingroup
\small
\setlength{\parindent}{0pt}%
\setlength{\parskip}{0pt}%
\setlength{\emergencystretch}{3em}%
\newcommand{\VerdierContentsRow}[3]{%
  \par\noindent
  \hangindent=#1\relax\hangafter=1
  \hspace*{#1}#2\nobreak
  \leaders\hbox{\kern.32em.\kern.32em}\hfill
  \nobreak\hbox{\kern.65em #3}\par}
\newcommand{\VerdierContentsChapter}[1]{%
  \par\addvspace{.55\baselineskip}\noindent\textbf{#1}\par\smallskip}
\newcommand{\VerdierContentsSection}[2]{%
  \par\addvspace{.45\baselineskip}\noindent
  \hfill\textbf{#1.}\enspace #2\hfill\mbox{}\par\smallskip}

\vspace*{4.25\baselineskip}
\VerdierSemanticAnchor{toc-title}%
\begin{center}
  {\Large\bfseries 目录}
\end{center}
\vspace{2.5\baselineskip}

\VerdierSemanticAnchor{toc-introduction}%
\VerdierContentsRow{0pt}{\textbf{引言}}{1}
\VerdierSemanticAnchor{toc-acknowledgments}%
\VerdierContentsRow{0pt}{\textbf{致谢}}{15}

\VerdierSemanticAnchor{toc-ch1}%
\VerdierContentsChapter{第一章：加性范畴中的复形范畴}
\VerdierSemanticAnchor{toc-ch1-sec1}%
\VerdierContentsSection{1}{分次范畴}
\VerdierSemanticAnchor{toc-ch1-sec1-1}%
\VerdierContentsRow{.5em}{\textbf{1.1.}\enspace $G$-分次范畴}{17}
\VerdierSemanticAnchor{toc-ch1-sec1-2}%
\VerdierContentsRow{.5em}{\textbf{1.2.}\enspace 群在范畴上的同构意义下作用}{19}
\VerdierSemanticAnchor{toc-ch1-sec1-3}%
\VerdierContentsRow{.5em}{\textbf{1.3.}\enspace 两种语言的等价}{22}
\VerdierSemanticAnchor{toc-ch1-sec1-4}%
\VerdierContentsRow{.5em}{\textbf{1.4.}\enspace 群的变换}{27}
\VerdierSemanticAnchor{toc-ch1-sec1-5}%
\VerdierContentsRow{.5em}{\textbf{1.5.}\enspace $\mathbb Z^n$-范畴}{34}
\VerdierSemanticAnchor{toc-ch1-sec1-6}%
\VerdierContentsRow{.5em}{\textbf{1.6.}\enspace Koszul 余循环与符号规则}{43}

\VerdierSemanticAnchor{toc-ch1-sec2}%
\VerdierContentsSection{2}{加法范畴中的复形}
\VerdierSemanticAnchor{toc-ch1-sec2-1}%
\VerdierContentsRow{.5em}{\textbf{2.1.}\enspace $n$-重复形}{49}
\VerdierSemanticAnchor{toc-ch1-sec2-2}%
\VerdierContentsRow{.5em}{\textbf{2.2.}\enspace 关联单复形}{50}
\VerdierSemanticAnchor{toc-ch1-sec2-3}%
\VerdierContentsRow{.5em}{\textbf{2.3.}\enspace 复形范畴上的分次结构}{54}
\VerdierSemanticAnchor{toc-ch1-sec2-4}%
\VerdierContentsRow{.5em}{\textbf{2.4.}\enspace 函子到复形的延拓}{56}
\VerdierSemanticAnchor{toc-ch1-sec2-5}%
\VerdierContentsRow{.5em}{\textbf{2.5.}\enspace 同伦}{60}

\VerdierSemanticAnchor{toc-ch1-sec3}%
\VerdierContentsSection{3}{特异三角}
\VerdierSemanticAnchor{toc-ch1-sec3-1}%
\VerdierContentsRow{.5em}{\textbf{3.1.}\enspace 态射的映射锥}{65}
\VerdierSemanticAnchor{toc-ch1-sec3-2}%
\VerdierContentsRow{.5em}{\textbf{3.2.}\enspace 映射锥函子的性质}{69}
\VerdierSemanticAnchor{toc-ch1-sec3-3}%
\VerdierContentsRow{.5em}{\textbf{3.3.}\enspace 特异三角的性质}{78}
\VerdierSemanticAnchor{toc-ch1-sec3-4}%
\VerdierContentsRow{.5em}{\textbf{3.4.}\enspace 特异三角的函子性质}{84}
\VerdierSemanticAnchor{toc-ch1-appendix}%
\VerdierContentsRow{0pt}{\textbf{附录：关于符号选择的说明}}{91}
\endgroup
